\documentclass[12pt, twoside]{article}
\input{../preamble}
\setlength{\parskip}{0.3\baselineskip}

\fancyhead[LE]{\thepage}
\fancyhead[RO]{Markscheme}
\fancyhead[LO]{La Scuola d'Italia / Dr.\ Huson / IB Math 12 \\* 9 September 2026}

\begin{document}
\subsubsection*{1.1 Classwork: Frequency Tables --- Markscheme}

\begin{enumerate}[itemsep=0.3cm]

%--- Problem 1: Tally and frequency table from categorical raw data [6 marks] ---
\item[1.] [Maximum mark: 6]

Raw data as printed (24 replies), with the counts:

{\small
\begin{center}
\setlength{\tabcolsep}{4pt}
\begin{tabular}{cccccccc}
bus & subway & walk & car & subway & walk & bike & subway \\
walk & bus & subway & subway & car & walk & subway & bus \\
subway & bike & walk & bus & car & subway & walk & bus \\
\end{tabular}
\end{center}}

\vspace{-0.2cm}

{\renewcommand{\arraystretch}{1.3}
\begin{center}
\begin{tabular}{|l|c|c|c|c|c|c|}
\hline
\textbf{How you travel} & walk & bus & subway & car & bike & \textbf{Total} \\ \hline
\textbf{Frequency} & 6 & 5 & 8 & 3 & 2 & \textbf{24} \\ \hline
\end{tabular}
\end{center}}

\textbf{(a)} Tally marks recorded for all five categories
\hfill \textbf{A1}

Frequencies $6,\; 5,\; 8,\; 3,\; 2$ all correct \hfill \textbf{A1}

\emph{The tally A1 is for a tally that is used and that agrees with the
candidate's own frequency column, however the strokes are grouped; the
five-bar gate is not required. The frequency A1 requires all five values
correct. Accept a completed frequency column with the tally column left
empty for the second A1 only.}

\textbf{(b)} Total $= 6 + 5 + 8 + 3 + 2 = 24$ \hfill \textbf{A1}

\emph{Accept 24 stated with no working, or read from the stem.}

\textbf{(c)} Mode is \emph{subway}. \hfill \textbf{A1}

\emph{Accept ``the subway'' or ``subway, 8 students''. Do not accept
``8'' alone --- the mode of a categorical variable is the category, not
its frequency.}

\emph{FT: award A1 for the category with the largest frequency in the
candidate's own table from (a).}

\textbf{(d)} $\dfrac{6}{24} = \dfrac{1}{4} = 25\%$ \hfill \textbf{A1}

\emph{Accept any equivalent fraction, decimal or percentage; simplification
is not required. FT: award A1 for the candidate's own walk frequency
divided by the candidate's own total.}

\textbf{(e)} It is categorical, because the replies are names of groups
(words), not numbers that can be measured or counted. \hfill \textbf{R1}

\emph{R1 is earned by any answer that identifies the variable as
categorical \emph{and} gives a reason of the right kind --- the data are
words / labels / groups rather than numerical measurements.
``Categorical'' with no reason earns nothing.}

\newpage

%--- Problem 2: Frequency table from discrete numerical raw data [6 marks] ---
\item[2.] [Maximum mark: 6]

Raw data as printed (25 replies), with the counts:

{\small
\begin{center}
\begin{tabular}{ccccc}
1 & 2 & 0 & 3 & 1 \\
1 & 2 & 4 & 1 & 2 \\
1 & 3 & 0 & 2 & 1 \\
1 & 2 & 3 & 1 & 2 \\
1 & 4 & 0 & 2 & 3 \\
\end{tabular}
\end{center}}

\vspace{-0.2cm}

{\renewcommand{\arraystretch}{1.3}
\begin{center}
\begin{tabular}{|l|c|c|c|c|c|c|}
\hline
\textbf{Number of siblings} & 0 & 1 & 2 & 3 & 4 & \textbf{Total} \\ \hline
\textbf{Frequency} & 3 & 9 & 7 & 4 & 2 & \textbf{25} \\ \hline
\end{tabular}
\end{center}}

\textbf{(a)} Tally marks recorded for all five values \hfill \textbf{A1}

Frequencies $3,\; 9,\; 7,\; 4,\; 2$ all correct \hfill \textbf{A1}

\emph{The tally A1 is for a tally that is used and that agrees with the
candidate's own frequency column, however the strokes are grouped. The
frequency A1 requires all five values correct. Accept a completed
frequency column with the tally column left empty for the second A1 only.}

\textbf{(b)} Total $= 3 + 9 + 7 + 4 + 2 = 25$ \hfill \textbf{A1}

\emph{Accept 25 stated with no working, or read from the stem.}

\textbf{(c)} Mode is $1$. \hfill \textbf{A1}

\emph{Accept ``1 sibling'' or ``one brother or sister''. Do not accept
``9'' alone --- that is the frequency, not the mode.}

\emph{FT: award A1 for the value with the largest frequency in the
candidate's own table from (a).}

\textbf{(d)} At least 2 means $2$, $3$ or $4$: $7 + 4 + 2 = 13$ students.
\hfill \textbf{A1}

\emph{A1 is for the value 13, reached by any correct route --- adding the
frequencies for 2, 3 and 4, counting directly from the raw list, or
subtracting from the total ($25 - 3 - 9 = 13$).}

\emph{A candidate who reads ``at least 2'' as ``more than 2'' obtains 6;
this earns no mark, since the whole mark is for the count. FT: award A1
for the sum of the candidate's own frequencies for $2$, $3$ and $4$.}

\textbf{(e)} It is discrete, because the number of brothers and sisters
is counted and can only take separate whole-number values --- you cannot
have $2.4$ siblings. \hfill \textbf{R1}

\emph{R1 is earned by any answer that says the variable is discrete
\emph{and} gives a reason of the right kind --- the values are counted,
are whole numbers, or cannot take values in between. ``Discrete'' with no
reason earns nothing.}

\end{enumerate}
\end{document}
